% English translation of checked_pages/esq_119_048.tex; source: cote 119, site 48, PDF page 49, typescript page 46.
% AI-translated by GPT-6 Astra (OpenAI, ChatGPT Pro mode).
\scanpage{119/48}
\mspage{46}

me vast and rich fields ripe for harvest, without
ever managing to communicate this vision, and the spark
that accompanies it, to a male (or female) student of mine, or to make it
lead to a common exploration of greater or lesser
duration. Each time, right up to the present day, after an exploration
of a few days or a few weeks, in which I discovered,
as a scout, riches unsuspected at the outset, the journey
was cut short when it became clear that I would be alone in
continuing it. Stronger interests then took precedence over
a journey that consequently appeared as a digression, or even
a scattering of effort, rather than an adventure pursued together.

One of these themes was that of plane polygons, centered
on the moduli varieties that can be associated with them. One of
the surprises here was the irruption of algebraic geometry into
a context that had seemed to me quite far removed from it. This kind of
surprise, linked to the ubiquity of geometry in geometry as
such, has moreover recurred on several occasions.\tnote{TN-048-01}{The source repeats ``géométrie'' without an added adjective and has the form ``répéré''. ``Recurred'' follows the repetition expressed by ``à plusieurs reprises''; the intended French verb remains uncertain, and the checked French wording is retained.}

Another theme was that of curves (in particular
circles) immersed in a surface, with particular attention
to the ``stable'' case in which the singular points are ordinary
double points (and also the more general case in which the
different branches at a point cross one another), often with
the additional assumption that the immersion is ``cellular'',
i.e. gives rise to a map. A variant of situations of
this type is that of immersions of a surface with nonempty boundary,
and first and foremost of a disk (which had been pointed out to me by
A. Campo some ten years ago). Beyond the question of
various combinatorial formulations of such situations, which
is scarcely more than an exercise in syntax, I have been interested
above all in a dynamic view of the possible configurations,
with\unsure{.}\tnote{TN-048-02}{The doubtful dot after ``avec'' is retained after ``with''.} the passage from one to another by continuous deformations, which
can be decomposed into composites of two types of \uline{elementary}
\uline{operations} and their inverses, namely the ``\uline{sweeping}'' of a branch
of a curve across a double point, and the \uline{erasure} or \uline{creation}

\margin{48}
